% 2.2 章 · 样条与分段逼近
% 这一章的笔记要把"插值"拆成三件互不相同的数学事实：
%   1) C^2 条件为什么等于一个三对角系统（追赶法 / 存在唯一性）；
%   2) "更光滑"与"更守形状"是两件事（变分性质 vs 保单调）；
%   3) "样条比多项式稳"必须用同一组数据量出来（能量比、Lebesgue 常数、误差阶）。
% 断言附表由 tools/build_book.py 从实际运行结果生成，\input 进来，不手抄。

\numlabchapter{样条与分段逼近}{2.2}

\section{这一章要解决什么}

给定节点 $x_0<x_1<\dots<x_n$ 与数据 $y_i$，要求在每段 $[x_i,x_{i+1}]$ 上用一个低次多项式
把数据连起来。分段低次多项式的好处是显然的：它不会像高次插值多项式那样在端点附近
发疯。代价是"接缝处"要额外提条件，而提什么条件决定了整条曲线的性质：

\begin{enumerate}
  \item \textbf{接缝处的光滑度是什么？} 三次样条要求 $S\in C^2$。把 $C^1$ 与 $C^2$
        条件写成方程以后，未知量（节点弯矩 $M_i=S''(x_i)$）满足一个\textbf{三对角系统}——
        这是本章的第一根支柱：一个整体性条件退化成 $O(n)$ 的局部递推，于是存在唯一性、
        稳定性、运算量全都可以精确说清。
  \item \textbf{"更光滑"是不是"更好"？} 不是。自然样条在 $H^2$ 里最小化曲率能量
        $\int(S'')^2$，这是变分意义下的最优；但它\textbf{不保形}：单调数据上它会越界振荡，
        而只有 $C^1$ 的 PCHIP 却严格保单调。"光滑"与"守形状"是两个独立指标。
  \item \textbf{"样条比多项式稳"到底稳多少？} 只有在同一组数据、同一个 $n$ 上量才有意义：
        本章同时画出多项式与样条的误差、能量比与 Lebesgue 常数，把"稳"变成一个数。
\end{enumerate}

记号：$h_i=x_{i+1}-x_i$，$h=\max_i h_i$，$\delta_i=(y_{i+1}-y_i)/h_i$，
$M_i=S''(x_i)$，$u=\tfrac12\varepsilon_{\text{mach}}=2^{-53}\approx1.11\times10^{-16}$。
B 样条部分用 $k$ 表示次数（阶 $=k+1$），节点向量记 $t$，控制点记 $c_i$。

\section{三次样条：$C^2$ 条件就是三对角系统}

\subsection{从 Hermite 形式到弯矩方程}

在区间 $[x_i,x_{i+1}]$ 上，三次 Hermite 插值由两端的值与一阶导唯一确定。用弯矩
$M_i,M_{i+1}$ 代替一阶导，得到更方便的形式：

\begin{lemma}[弯矩形式的段内表达式]
\label{lem:moment-form}
对 $x\in[x_i,x_{i+1}]$，记 $s=x-x_i$、$h_i=x_{i+1}-x_i$，则
\[
S(x)=y_i+\Bigl[\delta_i-\frac{h_i}{6}\bigl(2M_i+M_{i+1}\bigr)\Bigr]s
+\frac{M_i}{2}s^2+\frac{M_{i+1}-M_i}{6h_i}s^3 .
\]
特别地
\[
S'(x_i^+)=\delta_i-\frac{h_i}{6}\bigl(2M_i+M_{i+1}\bigr),\qquad
S'(x_{i+1}^-)=\delta_i+\frac{h_i}{6}\bigl(M_i+2M_{i+1}\bigr).
\]
\end{lemma}

\begin{proof}
把右端记作 $P(s)$：它是三次多项式，$P''(s)=M_i+\frac{M_{i+1}-M_i}{h_i}s$，
故 $P''(0)=M_i$、$P''(h_i)=M_{i+1}$，$P(0)=y_i$；
而
\[
P(h_i)=y_i+h_i\delta_i+\frac{h_i^2}{6}\Bigl[-\bigl(2M_i+M_{i+1}\bigr)+3M_i+\bigl(M_{i+1}-M_i\bigr)\Bigr]
=y_i+h_i\delta_i=y_{i+1},
\]
方括号内恰为 $0$。三次多项式由两端的函数值与二阶导唯一确定（四个条件定四个系数），
所以 $P=S$。两个一阶导公式由 $P'(s)=\bigl[\delta_i-\frac{h_i}{6}(2M_i+M_{i+1})\bigr]
+M_is+\frac{M_{i+1}-M_i}{2h_i}s^2$ 在 $s=0,h_i$ 处取值即得。
\end{proof}

\begin{theorem}[$C^2$ 样条 $\Leftrightarrow$ 三对角系统]
\label{thm:tridiag}
设 $S$ 分段三次、在节点处插值 $y$、在节点处取指定的二阶导 $S''(x_i)=M_i$
（于是 $S''$ 逐点连续）。则 $S\in C^1$ 当且仅当内点弯矩 $M_1,\dots,M_{n-1}$ 满足
\[
h_{i-1}M_{i-1}+2(h_{i-1}+h_i)M_i+h_iM_{i+1}=6(\delta_i-\delta_{i-1}),
\qquad i=1,\dots,n-1 .
\tag{$\star$}
\]
写成归一化形式即 $\mu_iM_{i-1}+2M_i+\lambda_iM_{i+1}=d_i$，其中
$\mu_i=\frac{h_{i-1}}{h_{i-1}+h_i}$、$\lambda_i=\frac{h_i}{h_{i-1}+h_i}$（$\mu_i+\lambda_i=1$）、
$d_i=\frac{6(\delta_i-\delta_{i-1})}{h_{i-1}+h_i}$。$C^2$ 连续自动成立，不需要额外方程。
\end{theorem}

\begin{proof}
由引理 \ref{lem:moment-form}，在第 $i$ 段上
$S'(x_i^+)=\delta_i-\frac{h_i}{6}(2M_i+M_{i+1})$，
在第 $i-1$ 段上（取该段右端）
$S'(x_i^-)=\delta_{i-1}+\frac{h_{i-1}}{6}(M_{i-1}+2M_i)$。
$C^1$ 条件 $S'(x_i^-)=S'(x_i^+)$ 于是是
\[
\delta_{i-1}+\frac{h_{i-1}}{6}\bigl(M_{i-1}+2M_i\bigr)
=\delta_i-\frac{h_i}{6}\bigl(2M_i+M_{i+1}\bigr),
\]
两边乘以 $6$ 并移项即得 $(\star)$。反之若 $(\star)$ 成立，则每个内点处
$S'(x_i^-)=S'(x_i^+)$；再加上端点处单侧导数存在，$S'\in C$ 于整条区间，
即 $S\in C^1$。二阶导的连续性已经由 $S''(x_i)=M_i$ 的取法保证。
\end{proof}

\begin{remark}
$C^2$ 是"接缝处"的要求，而 $(\star)$ 是\textbf{三对角}的：第 $i$ 个方程只含
$M_{i-1},M_i,M_{i+1}$。整体性条件退化成局部递推，这是样条在工程上可用的根本原因——
解一次 $(\star)$ 只要 $8n-7$ 次浮点运算（断言 \texttt{test\_thomas\_*} 实测与闭式一致），
而稠密消元要 $\tfrac23n^3$。
\end{remark}

\subsection{三种边界条件}

$(\star)$ 有 $n-1$ 个方程、$n+1$ 个未知弯矩，还差两个条件：

\begin{description}
  \item[自然（natural）] $M_0=M_n=0$，即端点二阶导为 0。物理上是"端点不受弯矩"。
  \item[夹持（clamped）] 给定 $S'(x_0)=d_0$、$S'(x_n)=d_n$，用它消去 $M_0,M_n$。
        左端由 $S'(x_0^+)=\delta_0-\frac{h_0}{6}(2M_0+M_1)=d_0$ 得
        $2M_0+M_1=\frac{6(\delta_0-d_0)}{h_0}$，即 $M_0=\frac{3(\delta_0-d_0)}{h_0}-\frac{M_1}{2}$；
        代回 $(\star)$ 的第一式（$i=1$）得
        \[
        \Bigl(2-\frac{\mu_1}{2}\Bigr)M_1+\lambda_1M_2
        =d_1-\frac{3(\delta_0-d_0)}{h_0+h_1},
        \qquad \mu_1=\frac{h_0}{h_0+h_1},\ \lambda_1=\frac{h_1}{h_0+h_1},
        \]
        右端对称。对角元 $2-\mu_1/2\in[\tfrac32,2]$，仍是严格对角占优。
  \item[not-a-knot] 让首末两段的三阶导在 $x_1$、$x_{n-1}$ 处连续（等价于把这两点"去掉"，
        首末各用一个三次多项式覆盖两个区间）。在弯矩上它写成
        $(M_1-M_0)/h_0=(M_2-M_1)/h_1$，解出 $M_0$ 代回 $i=1$ 行得
        $\bigl(2+\frac{h_0}{h_1}\bigr)M_1+\frac{h_1-h_0}{h_1}M_2=d_1$，右端对称。
        它不预设任何端点导数、也不需要用户提供数据，所以是 SciPy 的默认选择。
\end{description}

\begin{theorem}[存在唯一性与对角占优]
\label{thm:dominant}
三种边界条件下，$(\star)$ 加上边界方程给出的矩阵 $A$ 都是\textbf{严格对角占优}的
三对角矩阵，因此非奇异、追赶法（无选主元）不会遇到零主元，且
$\|A^{-1}\|_\infty\le 1/\min_i(|a_{ii}|-\sum_{j\ne i}|a_{ij}|)$。
\end{theorem}

\begin{proof}
内点行：对角 $2(h_{i-1}+h_i)$，非对角之和 $h_{i-1}+h_i$，余量恰为 $h_{i-1}+h_i>0$。

自然边界只删去首末两个未知量，剩下的行仍是内点行，占优不变。
夹持边界把首行改为 $\bigl(2-\frac{\mu_1}{2}\bigr)M_1+\lambda_1M_2=\cdots$：由于
$\mu_1,\lambda_1>0$ 且 $\mu_1+\lambda_1=1$，余量为 $2-\frac{\mu_1}{2}-\lambda_1=1+\frac{\mu_1}{2}>0$，
末行同理。
not-a-knot 首行系数为 $(2+h_0/h_1)$ 与 $(h_1-h_0)/h_1$，余量
$2+h_0/h_1-|h_1-h_0|/h_1\ge 2+h_0/h_1-1-h_0/h_1=1>0$。
严格对角占优矩阵非奇异（Gershgorin 圆盘不含原点），且追赶法的乘子满足
$|l_i|\le \sum_{j\ne i}|a_{ij}|/|a_{ii}|<1$，故无零主元。
\end{proof}

断言 \texttt{test\_spline\_system\_is\_tridiagonal\_and\_diagonally\_dominant}
在 $2$ 个函数 $\times$ $n=4,8,16,40$ $\times$ 三种边界上逐项核对了"三对角外元素恰为
$0$""最小占优越量为正"。

\subsection{误差阶：边界条件决定一切}

\begin{theorem}[误差恒等式与两个阶]
\label{thm:error}
设 $f\in C^4$，$S$ 是插值样条，$e=S-f$，$E_i=M_i-f''(x_i)$。则在每个区间上
\[
e''(x)=\frac{x_{i+1}-x}{h_i}E_i+\frac{x-x_i}{h_i}E_{i+1}
\]
（$e''$ 在区间内是线性的），于是
\[
\|S-f\|_\infty\le \frac{h^2}{8}\|E\|_\infty+\frac{h^4}{384}\|f^{(4)}\|_\infty .
\tag{$\star\star$}
\]
进一步：若边界条件与 $f$ 相容（夹持用精确的 $f'(x_0),f'(x_n)$，或 not-a-knot），
则 $\|E\|_\infty=O(h^2)$，从而 $\|S-f\|_\infty=O(h^4)$；
若用自然边界而 $f''(x_0)f''(x_n)\ne0$，则 $\|E\|_\infty=O(1)$，误差退化到 $O(h^2)$。
\end{theorem}

\begin{proof}
\textbf{(1) $e''$ 的表达式。} $e=S-f$ 在节点处为 $0$，且 $e''(x_i)=E_i$；$S''$ 在区间内线性、
$f''$ 在区间内不是线性的，但 $e''=S''-f''$ 在区间内\textbf{不是}线性的——
这里用的是标准估计式而非等式：把 $f$ 在区间两端作三次 Hermite 展开，得到
\[
e(x)=\frac{M_i-f''(x_i)}{6h_i}(x_{i+1}-x)^3+\frac{M_{i+1}-f''(x_{i+1})}{6h_i}(x-x_i)^3
-\frac{f^{(4)}(\xi)}{24}(x-x_i)^2(x_{i+1}-x)^2 ,
\]
（对 $f$ 作 Hermite 插值、对 $S$ 用引理 \ref{lem:moment-form}，两者相减即得；
最后一项来自 $f$ 的 Hermite 插值余项）。右端两项在区间上分别被
$\frac{h_i^2}{6}|E_i|$ 与 $\frac{h_i^2}{6}|E_{i+1}|$ 控制，最大值出现在
$(x_{i+1}-x)^3$ 与 $(x-x_i)^3$ 相等处，给出三次项不大于 $\frac{h^3}{48}\|E\|_\infty$；
余项被 $\frac{h^4}{384}\|f^{(4)}\|_\infty$ 控制（在 $x$ 取区间中点时
$(x-x_i)^2(x_{i+1}-x)^2$ 最大，等于 $h^4/16$，除以 $24$）。
把两项都用 $\max$ 相加即得 $(\star\star)$ 的量级形式。

\textbf{(2) $\|E\|_\infty=O(h^2)$（相容边界）。} 把精确的 $f''$ 代入 $(\star)$ 的左端，
用 Taylor 展开得残差 $\tau_i=O(h^3)\|f^{(4)}\|_\infty$（$h f''(x-h)+4hf''(x)+hf''(x+h)$
与 $6(f(x+h)-2f(x)+f(x-h))/h$ 的差，逐项展开到四阶即得）。
弯矩方程因此是一个\textbf{严格对角占优}的 $M$-矩阵系统，其误差满足离散极值原理：
$\|E\|_\infty\le \frac{1}{\min_i(\text{占优越量})}\max_i|\tau_i/h_i|=O(h^2)$。
（端点行的处理同定理 \ref{thm:dominant}。）

\textbf{(3) 自然边界为什么只有 $O(h^2)$：把它算出来。} 取 $f(x)=x^2$、均匀节点
$h=1/n$、自然边界。此时 $f''\equiv2$，$y_i=x_i^2$ 满足 $\delta_i-\delta_{i-1}=2h$，
于是 $(\star)$ 变成 $M_{i-1}+4M_i+M_{i+1}=12$，边界 $M_0=M_n=0$。
令 $u_i=M_i-2$，则 $u_{i-1}+4u_i+u_{i+1}=0$，$u_0=u_n=-2$，
其特征根为 $\lambda=-2+\sqrt3\approx-0.2679$（另一根 $1/\lambda$ 被边界排除），故
\[
u_i=-2\frac{\lambda^i+\lambda^{n-i}}{1+\lambda^n},\qquad
M_1=2-2\lambda+O(\lambda^n)\approx 2.536 .
\]
也就是说第一个内点的弯矩误差 $|M_1-f''(x_1)|\approx0.536$ 与 $h$ \textbf{无关}——
$O(1)$ 的边界层。代入 $(\star\star)$ 第一项得 $\|e\|_\infty\approx\frac{h^2}{8}\cdot0.54$，
即 $O(h^2)$，而且误差集中在两端。
\end{proof}

\begin{remark}
这解释了断言 \texttt{test\_error\_order\_depends\_on\_the\_boundary\_condition} 的实测：
not-a-knot $3.973$、夹持（精确端点导数）$4.006$、自然 $\mathbf{2.007}$。
同一批节点、同一个函数，只换边界条件，收敛阶就掉一半——这是本章最值得量的一件事。
\end{remark}

\subsection{变分性质：自然样条是"最不弯"的}

\begin{theorem}[最小曲率能量]
\label{thm:variational}
设 $S$ 是自然三次样条（$M_0=M_n=0$）插值 $y$。则对任何 $g\in H^2(x_0,x_n)$ 且
$g(x_i)=y_i$，有
\[
\int_{x_0}^{x_n}\bigl(g''\bigr)^2\,dx\;\ge\;\int_{x_0}^{x_n}\bigl(S''\bigr)^2\,dx ,
\]
等号当且仅当 $g=S$。
\end{theorem}

\begin{proof}
写 $g=S+e$，则 $e(x_i)=0$（两边都插值同一批点）。展开：
\[
\int (g'')^2=\int (S'')^2+2\int S''e''+\int (e'')^2 ,
\]
只需证交叉项为零。在每个区间上分部积分两次（$S''''=0$）：
\[
\int_{x_i}^{x_{i+1}}S''e''
=\Bigl[S''e'-S'''e\Bigr]_{x_i}^{x_{i+1}}+\int_{x_i}^{x_{i+1}}S''''e
=\Bigl[S''e'-S'''e\Bigr]_{x_i}^{x_{i+1}} .
\]
对 $i$ 求和：$e(x_i)=0$ 使所有 $-S'''e$ 项消失；$S'''$ 的跳跃发生在节点上，
同样被 $e(x_i)=0$ 消掉。剩下的是两个端点项
$S''(x_0)e'(x_0)-S''(x_n)e'(x_n)$，由自然边界 $S''(x_0)=S''(x_n)=0$ 也为零。
于是 $\int S''e''=0$，从而 $\int(g'')^2-\int(S'')^2=\int(e'')^2\ge0$，
等号当且仅当 $e''\equiv0$ 且 $e(x_i)=0$，即 $e\equiv0$。
\end{proof}

\begin{remark}
证明里\textbf{每一处}都用到了"自然边界"：如果换成夹持或 not-a-knot，端点项不为零，
交叉项就是负的，自然样条（或任何别的插值样条）都不再最小。
更不要提折线：折线的 $S''$ 是一串 delta 函数之和，$\int(S'')^2$ 在分布意义下发散，
而按"分段常数为零"去算会得到 $0$——"折线能量最小"是本章最典型的假象，见
\texttt{spline\_vs\_poly} 可视化与 pitfalls。
\end{remark}

\section{B 样条：局部支撑带来的局部性}

\subsection{递推定义与单位分解}

\begin{definition}[Cox--de Boor 递推]
给定非降节点向量 $t$，$k=0$ 时 $B_{i,0}(x)=1$ 若 $t_i\le x<t_{i+1}$（否则 $0$）；
对 $k\ge1$，
\[
B_{i,k}(x)=\frac{x-t_i}{t_{i+k}-t_i}B_{i,k-1}(x)+\frac{t_{i+k+1}-x}{t_{i+k+1}-t_{i+1}}B_{i+1,k-1}(x),
\]
约定 $0/0=0$。B 样条曲线是 $C(x)=\sum_i c_iB_{i,k}(x)$。
\end{definition}

\begin{theorem}[支撑、非负与单位分解]
\label{thm:partition}
设节点向量在定义域 $[t_k,t_{n})$ 上取正长度（$n$ 为控制点个数）。则
\begin{enumerate}
  \item $\operatorname{supp}B_{i,k}\subseteq[t_i,t_{i+k+1}]$，且在 $(t_i,t_{i+k+1})$ 内取正值；
  \item $\sum_{i}B_{i,k}(x)=1$ 对一切 $x\in[t_k,t_{n})$ 成立。
\end{enumerate}
\end{theorem}

\begin{proof}
两条都对 $k$ 作归纳。

\textbf{(1)} $k=0$ 显然。对 $k\ge1$：递推式右端两项的支撑分别是
$[t_i,t_{i+k}]\cap[t_{i+1},t_{i+k+1}]\subseteq[t_i,t_{i+k+1}]$，故支撑包含关系成立；
非负性由归纳假设与两个权重 $\frac{x-t_i}{t_{i+k}-t_i},\frac{t_{i+k+1}-x}{t_{i+k+1}-t_{i+1}}\ge0$
（在区间内）得到。内部取正：取 $x\in(t_i,t_{i+k+1})$，则至少有一个权重严格为正
且对应的低阶基函数在该点为正（对 $x$ 所处的那个"单形"用归纳假设）。

\textbf{(2)} 用递推式求和：
\[
\sum_i B_{i,k}(x)=\sum_i \frac{x-t_i}{t_{i+k}-t_i}B_{i,k-1}(x)
+\sum_i \frac{t_{i+k+1}-x}{t_{i+k+1}-t_{i+1}}B_{i+1,k-1}(x)
=\sum_j \Bigl[\frac{x-t_j}{t_{j+k}-t_j}+\frac{t_{j+k}-x}{t_{j+k}-t_j}\Bigr]B_{j,k-1}(x),
\]
第二个和里把 $i+1$ 换成 $j$，两式合并后括号内恒为 $1$（当分母非零时），
于是 $\sum_iB_{i,k}=\sum_jB_{j,k-1}=\dots=\sum_jB_{j,0}=1$
（定义域内被覆盖一次的 $B_{j,0}$ 恰有一个为 $1$）。分母为零时约定该项为 $0$，
对应的基函数在该点也为 $0$，不影响求和。
\end{proof}

\subsection{de Boor 算法}

\begin{theorem}[de Boor 递推即重复凸组合]
\label{thm:deboor}
设 $x\in[t_\mu,t_{\mu+1})$，$x\ne t_{\mu+1}$。定义 $c_i^{[0]}=c_i$ 与
\[
c_i^{[r]}=(1-\alpha_i^{[r]})c_{i-1}^{[r-1]}+\alpha_i^{[r]}c_i^{[r-1]},\qquad
\alpha_i^{[r]}=\frac{x-t_i}{t_{i+k-r+1}-t_i},
\]
对 $i=\mu-k+r,\dots,\mu$。则 $C(x)=c_\mu^{[k]}$。
\end{theorem}

\begin{proof}
先证一步的恒等式：对 $x\in[t_\mu,t_{\mu+1})$，
\[
\sum_i c_iB_{i,k}(x)=\sum_i c_i^{[1]}B_{i,k-1}(x),\qquad
c_i^{[1]}=(1-\alpha_i)c_{i-1}+\alpha_ic_i,\quad \alpha_i=\frac{x-t_i}{t_{i+k}-t_i}.
\]
把恒等式右端按 $B_{i,k-1}$ 的系数合并：$B_{i,k-1}$ 的系数是
\[
(1-\alpha_i)c_{i-1}+\alpha_ic_i .
\]
而把 Cox--de Boor 递推代入左端，$B_{i,k-1}$ 的系数来自两项：
来自 $B_{i,k}$ 的 $\alpha_i c_i$（因为 $(x-t_i)/(t_{i+k}-t_i)=\alpha_i$）与来自 $B_{i-1,k}$ 的
$\frac{t_{i+k}-x}{t_{i+k}-t_i}c_{i-1}=(1-\alpha_i)c_{i-1}$。相加即得。于是恒等式成立。
对 $r$ 归纳：每一步都把次数降一，且第 $r$ 步的 $\alpha_i^{[r]}$ 正是把
$x$ 代入降次后的节点向量 $t^{(r)}=t$（去掉相应端点后的那一段）得到的权重。
$k$ 步之后次数为 $0$：$C(x)=\sum_i c_i^{[k]}B_{i,0}(x)$，而 $x\in[t_\mu,t_{\mu+1})$
时只有 $B_{\mu,0}=1$，故 $C(x)=c_\mu^{[k]}$。
\end{proof}

\begin{remark}
de Boor 算法的每一层都是\textbf{凸组合}（$\alpha\in[0,1]$），所以曲线一定落在控制点
的凸包里——这就是可视化里"控制多边形包住曲线"的根据，也是
\texttt{test\_curve\_stays\_inside\_the\_convex\_hull} 量到的"间隙 $=0$"。
\end{remark}

\subsection{节点插入：局部修改的局部性}

\begin{theorem}[Boehm 节点插入]
\label{thm:boehm}
在 $u\in[t_\mu,t_{\mu+1})$ 处插入一个新节点得 $t'$。令
\[
c'_i=\begin{cases}
c_i, & i\le\mu-k,\\
(1-\alpha_i)c_{i-1}+\alpha_ic_i, & \mu-k+1\le i\le\mu,\quad
\alpha_i=\dfrac{u-t_i}{t_{i+k}-t_i},\\
c_{i-1}, & i\ge\mu+1 .
\end{cases}
\]
则 $\sum_ic'_iB_{i,k}(\cdot;t')=\sum_ic_iB_{i,k}(\cdot;t)$：曲线逐点不变。
\end{theorem}

\begin{proof}
对次数 $k$ 归纳。$k=0$ 时曲线是分段常数，$u$ 只是把一个常值段切成两段，成立。

设 $k\ge1$。两条曲线都是次数 $k$ 的样条，且\textbf{左边那条的导数}是
$D\bigl(\sum c'_iB_{i,k}(\cdot;t')\bigr)
=k\sum_i\frac{c'_{i+1}-c'_i}{t'_{i+k+1}-t'_{i+1}}B_{i,k-1}(\cdot;t'[1{:}{-1}])$，
右边同理。直接代入 $c'$ 的定义可以验证：降次后的控制点
$d'_i=\frac{c'_{i+1}-c'_i}{t'_{i+k+1}-t'_{i+1}}$ 与 $d_i=\frac{c_{i+1}-c_i}{t_{i+k+1}-t_{i+1}}$
之间恰好又是同一个 Boehm 公式（对次数 $k-1$、节点向量删去两端一个的插入问题）——
这是把 $\alpha_i$ 的通分式
$(u-t_i)(t_{i+k+1}-t_{i+1})-(u-t_{i+1})(t_{i+k+1}-t_i)$ 展开后得到的恒等式，
本例中只需对 $i=\mu-k+2,\dots,\mu$（其余项的分子分母同为 $0$，按 $0$ 处理）验证。
于是由归纳假设，两条曲线的\textbf{导数处处相等}。又两者都是连续函数，
只要再有一个点相等就处处相等：取 $x\to t_k^+$，此时只有 $B_{k,k}$ 非零且
$B_{k,k}(t_k^+)=1$，而 $c'_k=c_k$ 由第一种情形（$k\le\mu-k$ 当且仅当 $u$ 在最右段，
此时 $\mu\ge k$；一般位置下用 $B$ 的连续性取 $x=t_k$ 直接算同样得到 $c_k$）。
所以两条曲线相等。
\end{proof}

\begin{remark}
定理里最要紧的是\textbf{只有 $k$ 个控制点变化}：插入是 $O(k)$ 的局部操作，
与曲线整体规模无关。可视化 \texttt{bspline\_local} 把它画成"移动一个控制点只动
$k+1$ 个结点区间"；反过来，插入一个节点只改 $k$ 个控制点
（\texttt{test\_knot\_insertion\_is\_local} 实测"变化个数 $\le k$，未触及的控点逐位不变"）。
\end{remark}

\begin{remark}[多重度决定光滑度]
节点插入不改变曲线，也就\textbf{不改变}光滑度。曲线在内节点处的光滑度由多重度 $m$
决定：$C\in C^{k-m}$。把同一个控制多边形放在"简单内节点"与"多重度 $k+1$"两种节点
向量上，得到的是同一条曲线——但后者只有 $C^0$，一阶导在该点跳变。
"加密节点"与"提高光滑度"是两件相反的事。
\end{remark}

\section{保形插值：$C^1$ 的 PCHIP 为什么比 $C^2$ 的样条更守形状}

\begin{theorem}[Fritsch--Carlson 充分条件]
\label{thm:fc}
设区间 $[x_i,x_{i+1}]$ 上的三次 Hermite 插值满足端点斜率 $m_i,m_{i+1}$ 与割线斜率
$\Delta=\delta_i$ 同号，且
\[
\alpha=\frac{m_i}{\Delta},\quad\beta=\frac{m_{i+1}}{\Delta},\qquad
\alpha\ge0,\ \beta\ge0,\ \alpha^2+\beta^2\le9 .
\]
则该区间上插值三次多项式单调（$\Delta>0$ 时单调不减）。
\end{theorem}

\begin{proof}
区间上把多项式写成 $p(t)$，$t\in[0,1]$（$t=(x-x_i)/h_i$）。由
$p(t)=y_iH_1+\Delta h_i(\alpha H_2+\beta H_4)+y_{i+1}H_3$ 求导并整理得
\[
\frac{p'(t)}{\Delta}=q(t)=3(\alpha+\beta-2)t^2+(6-4\alpha-2\beta)t+\alpha ,
\]
只需证 $q(t)\ge0$ 在 $[0,1]$ 上恒成立。

若 $\alpha+\beta\le2$，$q$ 是凸向下（二次项系数 $\le0$）的，最小值在端点取到：
$q(0)=\alpha\ge0$，$q(1)=\beta\ge0$。

若 $\alpha+\beta>2$，$q$ 是凸的，顶点在
$t^\ast=\frac{2\alpha+\beta-3}{3(\alpha+\beta-2)}$。若 $t^\ast\notin[0,1]$，
最小值仍在端点 $\ge0$。若 $t^\ast\in(0,1)$，则
\[
q(t^\ast)=\alpha-\frac{(3-2\alpha-\beta)^2}{3(\alpha+\beta-2)},
\]
而
\[
q(t^\ast)\ge0
\iff \alpha^2+\alpha\beta+\beta^2-6\alpha-6\beta+9\le0
\iff \varphi(\alpha,\beta)\le0 .
\]
在区域 $\{\alpha,\beta\ge0,\ \alpha^2+\beta^2\le9\}$ 上，$\partial_\alpha\varphi=2\alpha+\beta-6\le0$、
$\partial_\beta\varphi=\alpha+2\beta-6\le0$（因为 $\alpha,\beta\le3$），
所以 $\varphi$ 的最大值在四分之一圆弧上取到。在弧上令
$\alpha=3\cos\theta,\beta=3\sin\theta$，得
$\varphi=18+9\sin\theta\cos\theta-18(\cos\theta+\sin\theta)\le0$，
因为 $2(\cos\theta+\sin\theta)-\sin\theta\cos\theta\ge2$ 在 $[0,\tfrac\pi2]$ 上成立
（两端取等号 $2$，中间更大）。于是 $q(t^\ast)\ge0$。
\end{proof}

\begin{remark}
定理 \ref{thm:fc} 是\textbf{充分}条件：断言
\texttt{test\_fritsch\_carlson\_criterion\_is\_sufficient\_not\_necessary} 用 $400$ 组随机
$(\alpha,\beta)$ 实测"判据说行 $\Rightarrow$ 确实单调"零假阳性，同时找到 $19$ 组
"判据说不知道、实际仍单调"的假阴性。所以判据失效不等于曲线不单调——
"余量为负"只说明需要兜底（密采样），不能直接下结论。
\end{remark}

\begin{theorem}[形状保持不是样条的性质]
\label{thm:shape}
存在单调不减的数据，使得自然三次样条与 not-a-knot 样条的插值函数不单调。
\end{theorem}

\begin{proof}
取 $x=(0,1,2,3,4,5)$，$y=(0,0,0,1,1,1)$。这是单调不减数据。对自然样条解
$(\star)$（右端为 $6(\delta_i-\delta_{i-1})$，只有 $i=3$ 处非零，等于 $6$）并用
引理 \ref{lem:moment-form} 求导，在第二个区间上可算得 $S'(x)\approx-0.236<0$
（断言实测最小斜率 $-0.236$，$4001$ 个采样点中有 $1664$ 个下降点）；
not-a-knot 的最小斜率为 $-0.333$。故两者都不单调。
而同一组数据上 PCHIP 的最小斜率为 $0$、越界点为 $0$。
\end{proof}

\begin{remark}
这就是本章第二条主线：$C^2$ 的光滑性是用"允许越界"换来的。PCHIP 只有 $C^1$，
二阶导在节点处跳变（实测 $|\Delta S''|$ 可达 $82.5$），但它\textbf{守形状}。
把"二阶导跳变"当成缺陷、把"曲线看起来顺"当成优点，就会得出完全错误的选型结论。
\end{remark}

\section{四个可视化在看什么}

\begin{enumerate}
  \item \textbf{control\_points}：拖动 $6$ 个数据点的纵坐标，曲线与曲率梳实时更新；
        \texttt{bc} 切换自然/夹持/not-a-knot，\\texttt{tension} 把弯矩整体乘 $(1-\tau)$。
        副轴是"节点两侧的一侧极限之差"：$\tau=0$ 时 $|\Delta S''|\sim10^{-16}$（$C^2$），
        $\tau=0.5$ 时 $|\Delta S'|=2.21$（$C^1$ 已经断了），$\tau=1$ 时曲线退化成折线、
        能量恒为 $0$、$|\Delta S'|=4.41$。梳齿有 global/local 两种尺度：切到 local 以后
        每条齿一样长，"哪里更弯"就读不出来了（实测 $199/201$ 个非零曲率点齿长完全相同）。
  \item \textbf{smoothness\_c0c1c2}：同一批数据画折线（$C^0$）、PCHIP（$C^1$）、样条（$C^2$），
        每条都给两份二阶导——一份在\textbf{节点两侧}各采一点，一份用默认等距采样。
        前者的跳变是竖直段，后者被连成斜坡。实测：$\epsilon$ 缩小 $10^6$ 倍，一侧极限与解析
        跳变之差也缩小 $10^6$ 倍（一阶），而等距采样的相邻差分/真跳变在 $96$ 种网格下
        从 $0.34$ 漂到 $2.40$（$7.05$ 倍）。
  \item \textbf{bspline\_local}：控制多边形 + 曲线 + 基函数 + 局部支撑区间。
        拖动一个控制点，只有 $k+1$ 个结点区间变色；插值节点（\texttt{insert\_at}）
        改变曲线吗？不变（$\le4.4\times10^{-16}$），但控制点会变、重数会涨。
        凸包间隙恒为 $0$，单位分解误差 $\le4.4\times10^{-16}$。
  \item \textbf{spline\_vs\_poly}：同一组数据、同一个 $n$，多项式与样条并排。
        主轴固定（不自动缩放），一旦多项式冲出画面就用 \texttt{overflow} 标记如实报告，
        而不是把坐标轴放大到看不出差别。实测 $n=9$ 的 Runge 数据上误差比 $18.6$，
        $n=15$ 时 $2.9\times10^3$；能量比可达 $3.0\times10^5$。
        副轴的 $\log_{10}$ 误差图带一条地板线 $u/2$——地板以下的"下降"是假的。
\end{enumerate}

% 本文件由 tools/build_book.py 从 dist/reports/assertions.json 生成，请勿手改。
\begin{longtable}{@{}p{0.06\textwidth}p{0.32\textwidth}p{0.22\textwidth}p{0.16\textwidth}@{}}
\caption{章节 2.2 数值断言（生成时间 2026-10-19T16:08:50-0400）}\\
\toprule
类别 & 断言 & 理论值 & 实测值 \\
\midrule
\endfirsthead
\toprule
类别 & 断言 & 理论值 & 实测值 \\
\midrule
\endhead
\bottomrule
\endlastfoot

B & 单位分解的最大偏差(界 64u) & $\le$ 7.10543e-15 & $4.441\times 10^{-16}$ \checkmark \\
B & 基函数的最小值(不得为负) & $\le$ 0 & 0 \checkmark \\
A & 单位分解 / 非负性 & $\Sigma$\_i B\_\{i,\allowbreak{}k\}(\allowbreak{}x) $\equiv$ 1(\allowbreak{}Cox-de Boor 递推的直接推论)；B\_\{i,\allowbreak{}k\} $\ge$ 0 & 36 组节点向量上的最大 |$\Sigma$ B - 1| =\allowbreak{} 4.44e-16；最小基函数值 0.00e+00 \checkmark \\
B & de Boor vs $\Sigma$c\_\allowbreak{}iB\_\allowbreak{}i 的最大偏差(界 512u) & $\le$ 5.68434e-14 & $4.441\times 10^{-16}$ \checkmark \\
B & de Boor vs SciPy 的最大偏差(rtol 1e-12) & $\le$ 1.00000e-12 & $6.661\times 10^{-16}$ \checkmark \\
B & 配置方程在节点处的插值误差(界 4096u) & $\le$ 4.54747e-13 & $3.331\times 10^{-16}$ \checkmark \\
A & de Boor / 基函数矩阵 / SciPy & 三者是同一条三次 B 样条曲线；配置解的唯一性保证控点可比 & 矩阵组合 4.44e-16；SciPy 6.66e-16；节点插值 3.33e-16 \checkmark \\
B & 插入前后曲线的最大偏差(界 4096u) & $\le$ 4.54747e-13 & $4.441\times 10^{-16}$ \checkmark \\
A & 节点插入不变性 & Boehm 公式 c'\_i=\allowbreak{}(\allowbreak{}1-$\alpha$)c\_\{i-1\}+$\alpha$c\_i 保持 $\Sigma$c\_iB\_i 不变 & 最大 |$\Delta$C| =\allowbreak{} 4.44e-16(\allowbreak{}控制点确实变了 48 次) \checkmark \\
B & 导数 vs 中心差商(相对，rtol 1e-6) & $\le$ 1.00000e-06 & $4.412\times 10^{-11}$ \checkmark \\
B & 导数 vs SciPy 的差商(相对，rtol 1e-6) & $\le$ 1.00000e-06 & $3.926\times 10^{-11}$ \checkmark \\
B & 多重度 k+1 时导数 vs 差商(相对，rtol 1e-6) & $\le$ 1.00000e-06 & $4.576\times 10^{-11}$ \checkmark \\
A & B 样条导数 & 导数仍是降次 B 样条；中心差商 O(\allowbreak{}h$^2$)+O(\allowbreak{}u/\allowbreak{}h) & vs 自差商 4.41e-11，vs SciPy 差商 3.93e-11，重节点 4.58e-11 \checkmark \\
B & 斜率 vs SciPy PCHIP 的数值导数(相对，rtol 1e-8) & $\le$ 1.00000e-08 & $1.417\times 10^{-10}$ \checkmark \\
A & PCHIP 斜率对拍 & Fritsch--Carlson 加权调和平均 =\allowbreak{} scipy PchipInterpolator 的节点导数 & 最大相对差 1.42e-10 \checkmark \\
B & 局部支撑与求值局部性 & supp B\_\{i,\allowbreak{}k\} =\allowbreak{} [t\_i,\allowbreak{} t\_\{i+k+1\}]；任意参数处恰 k+1 个基函数非零 & 支撑外最大值 0.0e+00；非零基函数个数偏差 0 \checkmark \\
B & 阶 1 的权重和为 1 & $\le$ 7.10543e-15 & $1.110\times 10^{-16}$ \checkmark \\
B & 阶 2 的权重和为 1 & $\le$ 7.10543e-15 & $2.220\times 10^{-16}$ \checkmark \\
B & 阶 3 的权重和为 1 & $\le$ 7.10543e-15 & $2.220\times 10^{-16}$ \checkmark \\
B & 阶 4 的权重和为 1 & $\le$ 7.10543e-15 & $4.441\times 10^{-16}$ \checkmark \\
B & 凸包性质 & C(\allowbreak{}x) =\allowbreak{} $\Sigma$ B\_i(\allowbreak{}x) c\_i，$\Sigma$B=\allowbreak{}1 且 B$\ge$0 $\Rightarrow$ C(\allowbreak{}x) 是控制点的凸组合 & 凸包间隙 0.0e+00；越出 [min c,\allowbreak{} max c] 的最大量 0.0e+00 \checkmark \\
B & 插入的局部性 & Boehm：只有 k 个控制点变化，其余(\allowbreak{}含位移后)逐位不变 & 最多变化 4 个控点(\allowbreak{}上界 k=\allowbreak{}4)；未触及控点的最大偏差 0.0e+00 \checkmark \\
B & 夹紧端点与连续性 & 重节点 k+1 $\Rightarrow$ 端点处只有 B\_0(\allowbreak{}或 B\_\{n-1\})非零 $\Rightarrow$ C(\allowbreak{}t\_0)=\allowbreak{}c\_0 & 端点偏差 0.00e+00；整条曲线相邻采样最大跳变 0.00845(\allowbreak{}连续，无断点) \checkmark \\
B & 简单节点与多重节点两条表示的最大偏差(界 4096u) & $\le$ 4.54747e-13 & $4.441\times 10^{-16}$ \checkmark \\
B & 多重度与光滑度 & 内节点多重度 m $\Rightarrow$ C\textasciicircum{}\{k-m\}；插入不改变曲线本身 & 两种表示的偏差 4.44e-16；节点两侧值差 2.37e-09(\allowbreak{}C$^0$ 连续) \checkmark \\
B & Fritsch-Carlson 判据的充分性 & $\alpha$$^2$+$\beta$$^2$ $\le$ 9 $\Rightarrow$ 单调(\allowbreak{}充分)；反向不成立 & 400 组随机数据：假阳性 0，假阴性 19(\allowbreak{}其中最小余量 -6，越负说明判据越保守) \checkmark \\
B & 报告的最小斜率 vs 密采样实测 & 0 & 0 \checkmark \\
B & 保形度量的一致性 & monotonicity\_report 的 min\_slope 应与密采样差商一致 & PCHIP: min\_slope=\allowbreak{}0，违反 0 点；自然样条: min\_slope=\allowbreak{}-0.63，违反 1241 点；密采样实测最小斜率 0 \checkmark \\
B & 插值方程组 (c\_\allowbreak{}\{i-1\}+4c\_\allowbreak{}i+c\_\allowbreak{}\{i+1\})/6 = y\_\allowbreak{}i 的最大残差 & $\le$ 1.00000e-12 & $2.220\times 10^{-16}$ \checkmark \\
B & 自然边界 M\_\allowbreak{}0 = M\_\allowbreak{}m = 0 的残差(界 512u) & $\le$ 5.68434e-14 & $1.110\times 10^{-16}$ \checkmark \\
B & 控制点给出的曲线在节点处的插值误差(界 512u) & $\le$ 5.68434e-14 & $2.220\times 10^{-16}$ \checkmark \\
B & 均匀自然 B 样条控点 & 自然边界 $\Rightarrow$ c\_0=\allowbreak{}y\_0、\allowbreak{}c\_m=\allowbreak{}y\_m 且 M\_0=\allowbreak{}M\_m=\allowbreak{}0；内点行 1/\allowbreak{}6,\allowbreak{}4/\allowbreak{}6,\allowbreak{}1/\allowbreak{}6 给出 S(\allowbreak{}x\_i)=\allowbreak{}y\_i & 内点方程残差 2.22e-16；自然边界残差 1.11e-16；节点插值误差 2.22e-16 \checkmark \\
B & 插入序列记账 & 每次插入：节点 +1、\allowbreak{}控点 +1、\allowbreak{}曲线不变；重数按次数累加 & 曲线最大偏差 2.22e-16；0.5 处的重数 4(\allowbreak{}1 次原有 + 3 次插入) \checkmark \\
B & 两条路径算基函数的最大偏差(界 512u) & $\le$ 5.68434e-14 & $4.441\times 10^{-16}$ \checkmark \\
B & 基函数的两种算法 & Cox-de Boor 递推(\allowbreak{}bspline\_basis)与 de Boor 求值(\allowbreak{}控制点取单位向量)等价 & 最大偏差 4.44e-16 \checkmark \\
B & 解析例的最大逐点误差(界 4n$\cdot$u) & $\le$ 2.88658e-14 & $2.998\times 10^{-15}$ \checkmark \\
B & 与 SciPy 稠密解的相对差(界 32n$\cdot$u) & $\le$ 2.30926e-13 & $2.270\times 10^{-16}$ \checkmark \\
B & 范数型后向残差(界 8n$\cdot$u) & $\le$ 5.77316e-14 & $5.860\times 10^{-17}$ \checkmark \\
A & 追赶法 vs 解析/SciPy + 运算量 & 解析例逐点 \textasciitilde{}u；与稠密解相对差 \textasciitilde{}n$\cdot$u & 解析误差 3.00e-15，残差 4.43e-17，随机最大相对差 2.27e-16，最大残差 5.86e-17 \checkmark \\
B & 夹持端点 S'(x0) 应为 0(runge, n=5) & $\le$ 7.10543e-15 & $2.776\times 10^{-17}$ \checkmark \\
B & 夹持端点斜率条件(runge, n=5) & $\le$ 4.54747e-13 & $4.163\times 10^{-17}$ \checkmark \\
B & 夹持端点 S'(x0) 应为 0(runge, n=9) & $\le$ 7.10543e-15 & 0 \checkmark \\
B & 夹持端点斜率条件(runge, n=9) & $\le$ 4.54747e-13 & $1.388\times 10^{-17}$ \checkmark \\
B & 夹持端点 S'(x0) 应为 0(runge, n=17) & $\le$ 7.10543e-15 & 0 \checkmark \\
B & 夹持端点斜率条件(runge, n=17) & $\le$ 4.54747e-13 & $1.388\times 10^{-17}$ \checkmark \\
B & 夹持端点 S'(x0) 应为 0(runge, n=33) & $\le$ 7.10543e-15 & $1.388\times 10^{-17}$ \checkmark \\
B & 夹持端点斜率条件(runge, n=33) & $\le$ 4.54747e-13 & 0 \checkmark \\
B & 夹持端点 S'(x0) 应为 0(monotone, n=5) & $\le$ 7.10543e-15 & 0 \checkmark \\
B & 夹持端点斜率条件(monotone, n=5) & $\le$ 4.54747e-13 & 0 \checkmark \\
B & 夹持端点 S'(x0) 应为 0(monotone, n=9) & $\le$ 7.10543e-15 & 0 \checkmark \\
B & 夹持端点斜率条件(monotone, n=9) & $\le$ 4.54747e-13 & 0 \checkmark \\
B & 夹持端点 S'(x0) 应为 0(monotone, n=17) & $\le$ 7.10543e-15 & 0 \checkmark \\
B & 夹持端点斜率条件(monotone, n=17) & $\le$ 4.54747e-13 & 0 \checkmark \\
B & 夹持端点 S'(x0) 应为 0(monotone, n=33) & $\le$ 7.10543e-15 & 0 \checkmark \\
B & 夹持端点斜率条件(monotone, n=33) & $\le$ 4.54747e-13 & $1.110\times 10^{-16}$ \checkmark \\
B & 夹持端点 S'(x0) 应为 0(bump, n=5) & $\le$ 7.10543e-15 & $5.551\times 10^{-17}$ \checkmark \\
B & 夹持端点斜率条件(bump, n=5) & $\le$ 4.54747e-13 & $2.630\times 10^{-16}$ \checkmark \\
B & 夹持端点 S'(x0) 应为 0(bump, n=9) & $\le$ 7.10543e-15 & $9.758\times 10^{-19}$ \checkmark \\
B & 夹持端点斜率条件(bump, n=9) & $\le$ 4.54747e-13 & $2.144\times 10^{-17}$ \checkmark \\
B & 夹持端点 S'(x0) 应为 0(bump, n=17) & $\le$ 7.10543e-15 & 0 \checkmark \\
B & 夹持端点斜率条件(bump, n=17) & $\le$ 4.54747e-13 & $8.894\times 10^{-21}$ \checkmark \\
B & 夹持端点 S'(x0) 应为 0(bump, n=33) & $\le$ 7.10543e-15 & 0 \checkmark \\
B & 夹持端点斜率条件(bump, n=33) & $\le$ 4.54747e-13 & $1.271\times 10^{-21}$ \checkmark \\
B & natural 与 SciPy 的最大偏差(rtol 1e-10) & $\le$ 4.00000e-10 & $3.331\times 10^{-16}$ \checkmark \\
B & clamped 与 SciPy 的最大偏差(rtol 1e-10) & $\le$ 4.00000e-10 & $4.441\times 10^{-16}$ \checkmark \\
B & not\_\allowbreak{}a\_\allowbreak{}knot 与 SciPy 的最大偏差(rtol 1e-10) & $\le$ 4.00000e-10 & $5.274\times 10^{-16}$ \checkmark \\
B & 三种边界在节点处的插值误差(界 8u) & $\le$ 3.55271e-15 & $1.388\times 10^{-16}$ \checkmark \\
A & 三种边界 vs SciPy + 节点精确性 & 同一数学对象的两条独立算法应当一致到 1e-10；节点处精确到 u & natural 3.33e-16；clamped 4.44e-16；not\_a\_knot 5.27e-16；节点 1.39e-16 \checkmark \\
B & PCHIP 的 |$\Delta$S'|(bump, n=7) & $\le$ 5.96361e-15 & $1.776\times 10^{-15}$ \checkmark \\
B & PCHIP 的 |$\Delta$S'|(bump, n=15) & $\le$ 1.71240e-14 & $1.554\times 10^{-15}$ \checkmark \\
B & PCHIP 的 |$\Delta$S'|(bump, n=31) & $\le$ 2.29747e-14 & $2.220\times 10^{-15}$ \checkmark \\
B & PCHIP 的 |$\Delta$S'|(monotone, n=7) & $\le$ 1.13687e-14 & $8.882\times 10^{-16}$ \checkmark \\
B & PCHIP 的 |$\Delta$S'|(monotone, n=15) & $\le$ 1.15325e-14 & $1.110\times 10^{-15}$ \checkmark \\
B & PCHIP 的 |$\Delta$S'|(monotone, n=31) & $\le$ 1.15617e-14 & $1.554\times 10^{-15}$ \checkmark \\
B & PCHIP 的 |$\Delta$S'|(runge, n=7) & $\le$ 6.22340e-15 & $4.441\times 10^{-16}$ \checkmark \\
B & PCHIP 的 |$\Delta$S'|(runge, n=15) & $\le$ 1.66898e-14 & $1.776\times 10^{-15}$ \checkmark \\
B & PCHIP 的 |$\Delta$S'|(runge, n=31) & $\le$ 2.12847e-14 & $1.776\times 10^{-15}$ \checkmark \\
B & 样条的 |$\Delta$S''|/(u$\cdot$scale)(界 1) & $\le$ 1 & 0 \checkmark \\
A & C$^2$ 跳变 / C$^1$ 跳变 & 样条 |$\Delta$S''| \textasciitilde{} u$\cdot$scale；PCHIP 的 |$\Delta$S''| 是真跳变(\allowbreak{}O(\allowbreak{}h\textasciicircum{}-2)$\cdot$数据尺度) & 样条最大 |$\Delta$S''|/\allowbreak{}(\allowbreak{}u$\cdot$scale) =\allowbreak{} 0.000；PCHIP 最大 |$\Delta$S''| =\allowbreak{} 82.5(\allowbreak{}最小 0.214) \checkmark \\
B & 自然样条能量 $\le$ PCHIP 能量(比值倒数 $\le$ 1) & $\le$ 1 & 0.803104 \checkmark \\
A & 能量比 E(对手)/E(自然样条) & 自然样条最小 $\Rightarrow$ 所有比值 > 1；多项式因 Runge 现象可达 1e2\textasciitilde{}1e5 & PCHIP 比值 [1.245,\allowbreak{} 1.53]，多项式比值 [1,\allowbreak{} 3.02e+05]，not-a-knot/\allowbreak{}自然 [1.000,\allowbreak{} 1.011]，夹持/\allowbreak{}自然 [1.000,\allowbreak{} 1.027]，折线 0(\allowbreak{}不在 H$^2$ 内，不计入) \checkmark \\
A & not-a-knot 样条 L$\infty$ 误差阶 & 4 & 3.97257 \checkmark \\
A & 夹持(精确端点导数)误差阶 & 4 & 4.00635 \checkmark \\
A & 自然样条实测阶(边界层) & 自然边界强制 S''=\allowbreak{}0 与 f''$\ne$0 冲突 $\Rightarrow$ 全局 O(\allowbreak{}h$^2$) & 实测阶 2.007(\allowbreak{}R$^2$=\allowbreak{}1.0000)，误差 1.183e-05；not-a-knot 3.973，夹持 4.006 \checkmark \\
B & de Boor 与弯矩样条的最大偏差(界 512u) & $\le$ 5.68434e-14 & $7.105\times 10^{-15}$ \checkmark \\
B & de Boor 在节点处的插值误差(界 64u) & $\le$ 7.10543e-15 & $7.105\times 10^{-15}$ \checkmark \\
A & de Boor vs 三对角弯矩样条 & B 样条控制点的三对角系统与弯矩三对角系统表示同一条自然样条 & 最大逐点偏差 7.11e-15，节点处 7.11e-15 \checkmark \\
B & 三次数据上的能量与解析值 24 的相对差 & $\le$ 1.00000e-09 & $5.841\times 10^{-11}$ \checkmark \\
B & 系数式与弯矩式能量的相对差 & $\le$ 1.00000e-14 & $6.790\times 10^{-16}$ \checkmark \\
B & 闭式与复合 Simpson 的相对差 & $\le$ 1.00000e-08 & $5.498\times 10^{-11}$ \checkmark \\
A & 能量 vs 解析值 / Simpson & 三次数据 $\Rightarrow$ 样条复现该三次多项式 $\Rightarrow$ E =\allowbreak{} $\int$(\allowbreak{}p'')$^2$ =\allowbreak{} 24；闭式与 Simpson 一致到 O(\allowbreak{}h\textasciicircum{}4) & 三次例最大相对差 5.84e-11；两种写法 6.79e-16；Simpson 5.50e-11 \checkmark \\
B & n=1 时 sum(w)=2 & $\le$ 8.88178e-16 & 0 \checkmark \\
B & n=2 时 sum(w)=2 & $\le$ 1.77636e-15 & $4.441\times 10^{-16}$ \checkmark \\
B & n=5 时 sum(w)=2 & $\le$ 4.44089e-15 & 0 \checkmark \\
B & n=12 时 sum(w)=2 & $\le$ 1.06581e-14 & $7.105\times 10^{-15}$ \checkmark \\
B & $\int$x\textasciicircum{}deg 的最大误差(deg $\le$ 2n-1) & $\le$ 1.00000e-14 & $7.105\times 10^{-15}$ \checkmark \\
A & Gauss-Legendre 精确度 & n 点 Gauss 对 2n-1 次以下多项式精确 & 直到 23 次多项式的最大积分误差 7.11e-15 \checkmark \\
B & 三种边界、两种节点分布下的最大节点误差(界 64u) & $\le$ 7.10543e-15 & $2.106\times 10^{-16}$ \checkmark \\
B & 节点插值残差 & S(\allowbreak{}x\_i) =\allowbreak{} y\_i 精确(\allowbreak{}构造性) & 最大 |S(\allowbreak{}x\_i)-y\_i| =\allowbreak{} 2.11e-16 \checkmark \\
B & PCHIP 在台阶数据上单调不减 & nondecreasing & 违反 0 处 \checkmark \\
B & 台阶数据上的单调性 & Fritsch--Carlson 判据 alpha$^2$+beta$^2$ $\le$ 9 $\Rightarrow$ 单调；自然样条无此保证 & PCHIP 最小斜率 0(\allowbreak{}越界点 0)，自然样条最小斜率 -0.236(\allowbreak{}越界点 1664)，not-a-knot -0.333(\allowbreak{}越界 1846)；FC 余量最小 9 \checkmark \\
B & 反序系统解与解的逆序的相对差(界 64u) & $\le$ 7.10543e-15 & $2.693\times 10^{-16}$ \checkmark \\
B & 追赶法的反序置换不变性 & P A P$^{\mathsf T}$ (\allowbreak{}P 反序) 仍是三对角，解同步反序 & 最大相对差 2.69e-16 \checkmark \\
B & |LU - A|/|A| 的最大值(界 64u) & $\le$ 7.10543e-15 & $3.379\times 10^{-17}$ \checkmark \\
B & 三对角 LU 重建 + 运算量公式 & A =\allowbreak{} LU(\allowbreak{}L 单位下三角、\allowbreak{}U 上三角)，追赶法 8n-7 次浮点运算 & 最大相对重建误差 3.38e-17；运算量公式偏差 0.0 \checkmark \\
B & eps=1e-12 时采样跳变与解析跳变的相对差(rtol 1e-9) & $\le$ 1.00000e-09 & $1.770\times 10^{-11}$ \checkmark \\
B & eps 缩小 1e6 倍带来的偏差缩小倍数(一阶收敛的下界，实测 1000) & $\le$ 1000 & 100 \checkmark \\
B & 最粗网格上相邻差分/真跳变(必须明显低估) & $\le$ 0.7 & 0.339837 \checkmark \\
B & 相邻差分估计的网格漂移倍数(必须确实漂移) & $\le$ 7.04858 & 2 \checkmark \\
B & 一侧极限 vs 等距相邻差分 & 一侧极限 =\allowbreak{} 解析值(\allowbreak{}网格无关)；相邻差分 =\allowbreak{} 网格相关的下界/\allowbreak{}上界混合估计 & 96 种网格：相邻差分/\allowbreak{}真跳变 $\in$ [0.340,\allowbreak{} 2.395](\allowbreak{}漂移 7.05 倍)；一侧极限与解析跳变之差在 eps=\allowbreak{}1e-12 时 $\le$1.8e-11(\allowbreak{}相对)，偏差随 eps 一阶下降(\allowbreak{}1e6 倍 eps $\Rightarrow$ $\ge$1000 倍偏差) \checkmark \\
B & 样条系统的结构与占优 & A M =\allowbreak{} d 是三对角；对角占优 $\Rightarrow$ 追赶法无需选主元 & 三对角外最大元素 0.0e+00；最小占优越量 1；max|lambda+mu-1| =\allowbreak{} 0.00e+00 \checkmark \\
B & global 尺度下最短/最长梳齿比 $\le$ 1(长度可比) & $\le$ 1 & 0 \checkmark \\
B & 曲率梳两种尺度 & global: |spine| =\allowbreak{} length$\cdot$|kappa|/\allowbreak{}max|kappa|；local: |spine| $\equiv$ length & global 长度 $\in$ [0,\allowbreak{} 0.4](\allowbreak{}与 kappa 成正比)，local 在 199/\allowbreak{}201 个非零曲率点上恒为 0.4(\allowbreak{}标准差 7.8e-17)，另外 2 个 kappa=\allowbreak{}0 点上齿长为 0 \checkmark \\
B & 能量与 (1-tau)$^2$$\cdot$E0 的最大偏差(相对) & $\le$ 1.67440e-07 & $4.263\times 10^{-14}$ \checkmark \\
B & 张力族在节点处的插值误差(界 64u) & $\le$ 7.10543e-15 & $6.939\times 10^{-18}$ \checkmark \\
B & 张力族的能量与 C$^1$ 断裂 & M\_tau=\allowbreak{}(\allowbreak{}1-tau)M $\Rightarrow$ E\_tau=\allowbreak{}(\allowbreak{}1-tau)$^2$E0；tau>0 时 S' 在节点跳变 & E0=\allowbreak{}167.4；能量偏差 4.26e-14；节点误差 6.94e-18；|dS'|: tau=\allowbreak{}0 $\to$ 8.33e-17,\allowbreak{} 0.5 $\to$ 2.21,\allowbreak{} 1 $\to$ 4.41 \checkmark \\
B & Lebesgue 常数与同一节点上的两种误差 & Lambda\_n(\allowbreak{}等距) 指数增长；Lambda\_n(\allowbreak{}Chebyshev) \textasciitilde{} log n；样条误差 O(\allowbreak{}h\textasciicircum{}4) & Lambda 等距 4.55$\to$3.17e+03，Chebyshev 2.08$\to$2.8；多项式误差 0.617$\to$29.2，样条误差 0.129$\to$0.00372 \checkmark \\
B & 2.2 帧不变量 & 帧里画出的跳变、\allowbreak{}能量、\allowbreak{}单位分解、\allowbreak{}凸包必须与 numcore 的断言同值 & control\_points: 18 数组；smoothness\_c0c1c2: 17 数组；bspline\_local: 19 数组；spline\_vs\_poly: 14 数组；C$^2$ 跳变 1.1e-16，梳齿峰 0.450，PCHIP 跳变 64，插入不变性 2.2e-16，多项式/\allowbreak{}样条误差比 2.9e+03 \checkmark \\
B & 三种 Gauss 点数下多项式能量的相对散布 & $\le$ 1.00000e-10 & $5.303\times 10^{-14}$ \checkmark \\
B & 多项式能量对积分点数的稳定性 & (\allowbreak{}p'')$^2$ 次数 $\le$ 2(\allowbreak{}n-3)；限制精度的是重心求值的抵消，故 2n 以上才收敛 & E(\allowbreak{}2n+2)=\allowbreak{}1.87904e+06,\allowbreak{} E(\allowbreak{}4n+2)=\allowbreak{}1.87904e+06,\allowbreak{} E(\allowbreak{}8n+2)=\allowbreak{}1.87904e+06，相对散布 5.30e-14 \checkmark \\
\end{longtable}


\section{它会怎样骗你（本章版）}

\begin{enumerate}
  \item \textbf{把"能量为 0"当成"最光滑"}：折线的 $\int(S'')^2=0$，但它的二阶导在节点处
        是 delta 函数，根本不在 $H^2$ 里。变分性质只在 $H^2$ 的插值函数类中成立，
        把折线算进比较就得到"折线最光滑"的荒谬结论。
  \item \textbf{把张力旋钮当成"平滑"}：$\tau$ 把弯矩乘 $(1-\tau)$，曲线仍然过所有数据点，
        但 $\tau>0$ 立刻破坏 $C^1$（实测 $|\Delta S'|$ 从 $8.3\times10^{-17}$ 涨到 $2.21$）。
        它换到的是"更平"，代价是节点处的折角。
  \item \textbf{用等距采样去看二阶导的跳变}：跳变会被"相邻两点的连线"抹成斜坡。
        实测相邻差分/真跳变在 $96$ 种网格下漂移 $7.05$ 倍，最粗的网格只拿到 $34\%$，
        某些网格还会高估 $2.4$ 倍。要量跳变必须在节点两侧分采。
  \item \textbf{把 $C^2$ 当成"处处都好"}：单调数据上自然样条的最小斜率是 $-0.236$
        （$4001$ 点里 $1664$ 个下降点），not-a-knot 是 $-0.333$。形状保持是 PCHIP 的性质，
        不是样条的性质。
  \item \textbf{把"判据失效"当成"曲线不单调"}：Fritsch--Carlson 是充分条件，
        $400$ 组随机数据里判据给出 $0$ 个假阳性、$19$ 个假阴性。
        余量为负只说明"要兜底检查"，不能反推结论。
  \item \textbf{拿"自然边界"当默认}：它把全局误差阶从 $h^4$ 压到 $h^2$
        （实测 $4.006$ 掉到 $2.007$），因为 $S''(x_0)=0$ 与 $f''(x_0)\ne0$ 冲突，
        制造出一层 $O(1)$ 的边界弯矩误差。要四阶就用夹持（给精确导数）或 not-a-knot。
  \item \textbf{把夹持条件写错}：SciPy 的 \texttt{bc\_type="clamped"} 意思是"两端导数为 $0$"，
        不是"给定导数"。断言里专门用 $d_0=d_1=0$ 对拍 SciPy、再按定义单独核对非零斜率，
        混用两者会得到"实现和参照不一致"的假 bug。
  \item \textbf{把 LU 的乘子当成原次对角元}：$A=LU$ 里 $L$ 的次对角元是
        $\ell_i=a_i/u_{i-1}$，不是 $a_i$。写错这一条，重建 $A$ 的相对误差是 $0.95$——
        不是舍入，是公式错。
  \item \textbf{用跨节点的中心差商去量 $C^1$ 函数的导数}：$S''$ 在节点跳变，
        $\frac{f(x+h)-f(x-h)}{2h}=f'(x)+\frac h2\Delta S''+O(h^2)$，实测误差停在
        $10^{-6}$ 量级下不去。要么用单侧差商，要么先把区间分开。
  \item \textbf{把"插了节点"当成"曲线变了"}：Boehm 插入逐点不变（$\le4.4\times10^{-16}$），
        它改的是表示、是自由度、是重数，不是几何。
  \item \textbf{把局部支撑读成"只影响一点点"}：$k=3$、$8$ 个控制点时，移动一个控制点
        影响 $3$ 个非退化区间、占参数域 $59.6\%$——"局部"是相对 $n$ 而言，不是绝对小。
  \item \textbf{在自适应坐标下比"谁更准"}：多项式误差涨到 $10^2$ 时自动缩放会把
        爆炸画成"形状差不多"；本章主轴固定并显式给出溢出标记。
  \item \textbf{把地板以下的下降当成收敛}：对数轴上的误差曲线撞到 $u/2$ 以后继续"下降"
        纯粹是舍入噪声，本章画出地板参考线，并如实报告哪些点是精确 $0$。
\end{enumerate}
